% q2.tex — 五、问题二的模型建立与求解（2026-08-15, 素材见 PAPER_HANDOFF.md §6）

\subsection{5.1 数据准备}

问题二以品类为单位建模，品类级数据的口径如下。品类日销量为该品类各单品日销量的
合计；品类批发价以当日各单品销量为权重对批发价加权平均；品类成本加成定价由"单位
成本$\times(1+\text{加成率})$"定义，其中单位成本与加成率均只用非打折交易计算，以
排除促销价格对正常定价口径的干扰。损耗率取品类内单品损耗率的均值 $\bar\ell_c$。
此外，用历史打折交易均价与正常交易均价的比值定义品类打折数 $d_c$（即打折品实际
售价约为正常售价的 $d_c$ 倍），六品类的 $d_c$ 分别为 0.740、0.785、0.847、0.873、
0.771 与 0.836，说明蔬菜打折清仓的折扣水平整体稳定。

\subsection{5.2 销售总量与成本加成定价的关系}

以各品类日销量为因变量、当日成本加成定价为自变量，对 6--7 月的日度观测做曲线
拟合。对线性、二次、三次三种候选形式逐一比较拟合优度，三次多项式的整体表现最好，
故采用三次多项式刻画两者的关系：

\begin{equation}
B_c(P)=a_{c,0}+a_{c,1}P+a_{c,2}P^2+a_{c,3}P^3.
\label{eq:cubic}
\end{equation}

各品类的拟合系数与拟合优度见表2，拟合曲线见图10。

\begin{table}[H]
\centering
\caption{各品类销量-成本加成定价三次拟合结果}
\label{tab:cubic}
\small
\begin{tabular}{lllllc}
\toprule
品类 & $a_{c,3}$ & $a_{c,2}$ & $a_{c,1}$ & $a_{c,0}$ & $R^2$ \\
\midrule
花叶类 & $-11.92$ & $223.82$ & $-1352.63$ & $2797.83$ & 0.31 \\
花菜类 & $-0.04$ & $0.58$ & $-0.04$ & $26.56$ & 0.22 \\
水生根茎类 & $-0.06$ & $3.34$ & $-58.02$ & $344.11$ & 0.58 \\
茄类 & $0.05$ & $-1.76$ & $14.53$ & $-3.33$ & 0.03 \\
辣椒类 & $0.007$ & $0.70$ & $-14.47$ & $127.70$ & 0.04 \\
食用菌 & $-0.12$ & $4.40$ & $-48.78$ & $211.48$ & 0.24 \\
\bottomrule
\end{tabular}
\end{table}

\begin{figure}[H]
\centering
\includegraphics[width=0.8\textwidth]{figures/q2_fit_cubic.png}
\caption{各品类销售量与成本加成定价三次多项式拟合}
\label{fig:q2_fit}
\end{figure}

从拟合结果看，花叶类的需求体量最大：在其历史定价区间（约 5--6 元/kg）内，三次
曲线对应的日销量约 230 kg，且定价小幅波动时销量变化不大，说明叶菜类需求刚性较强，
属于生活必需菜品。水生根茎类的拟合优度最高（$R^2=0.58$），价格对销量的调节关系
较为明确；茄类与辣椒类的 $R^2$ 低于 0.1，两者的销量主要由季节与供给因素决定，
价格解释力有限，这一局限在第九章模型评价中进一步说明。三次关系式一方面回答了
"销售总量与成本加成定价关系"的分析要求，另一方面为 5.4 节优化模型提供"给定价格
下销量"的约束依据（方案 B 中用于计算各价格档位对应的需求量）。

\subsection{5.3 未来一周销量与批发价的预测}

问题二要求给出 2023 年 7 月 1 日至 7 日的补货与定价策略，而决策当周的销量与批发价
在建模时点均未知，需要先给出预测。4.2.4 节的统计表明销量与批发价序列均存在明显的
周内规律：周末销量显著高于工作日。据此，未来某一天的预测值取历史 6--7 月全部
同星期观测的均值：

\begin{equation}
\hat B_{c,t}=\frac{1}{n_{\mathrm{wd}(t)}}\sum_{s\in\mathcal{S}_{\mathrm{wd}(t)}}B_{c,s},
\qquad
\hat w_{c,t}=\frac{1}{n_{\mathrm{wd}(t)}}\sum_{s\in\mathcal{S}_{\mathrm{wd}(t)}}w_{c,s},
\label{eq:weekday_forecast}
\end{equation}

其中 $\mathrm{wd}(t)$ 为第 $t$ 天的星期，$\mathcal{S}_{\mathrm{wd}(t)}$ 为历史
6--7 月中该星期的全部日期。选取 6--7 月作为参照期的理由是：预测时段处于夏季，
参照期内蔬菜供应结构与需求水平与预测时段一致，避免了跨季节样本的干扰。六品类的
预测销量见表3，未来一周预测销量合计约 2,590 kg。

\begin{table}[H]
\centering
\caption{未来一周各品类预测日销量（kg）}
\label{tab:Bforecast}
\small
\begin{tabular}{lccccccc}
\toprule
品类 & 7月1日 & 7月2日 & 7月3日 & 7月4日 & 7月5日 & 7月6日 & 7月7日 \\
\midrule
花叶类 & 210.2 & 213.0 & 144.7 & 145.5 & 149.2 & 152.3 & 160.7 \\
花菜类 & 47.9 & 47.3 & 34.7 & 36.7 & 35.6 & 32.8 & 32.5 \\
水生根茎类 & 27.5 & 32.3 & 17.7 & 18.9 & 18.9 & 17.7 & 18.3 \\
茄类 & 38.7 & 41.2 & 26.9 & 24.4 & 25.3 & 26.5 & 27.0 \\
辣椒类 & 86.1 & 82.2 & 56.0 & 54.7 & 61.8 & 61.8 & 59.9 \\
食用菌 & 59.5 & 59.1 & 40.3 & 37.4 & 41.4 & 40.7 & 44.4 \\
\bottomrule
\end{tabular}
\end{table}

\subsection{5.4 补货与定价优化模型}

\subsubsection{5.4.1 目标函数}

商超一日的利润由三部分构成：正常售价部分的销售收入、损耗品按折扣清仓的回收收入，
以及进货成本。设品类 $c$ 第 $t$ 天补货量为 $x_{c,t}$，成本加成定价为 $P_{c,t}$，
有效成本为 $\hat w_{c,t}=w_{c,t}(1+\bar\ell_c)$，则七日总利润目标为

\begin{equation}
\max\ Z=\sum_{t=1}^{7}\sum_{c=1}^{6}x_{c,t}\left[P_{c,t}\bigl((1-\bar\ell_c)+\bar\ell_c d_c\bigr)-\hat w_{c,t}\right],
\label{eq:obj}
\end{equation}

式中 $(1-\bar\ell_c)P_{c,t}$ 为正常销售部分的单位收入，$\bar\ell_c d_c P_{c,t}$
为损耗部分按折扣 $d_c$ 清仓的单位回收，二者之和减去有效成本即为单位利润。与只
考虑正常销售的利润函数相比，式\eqref{eq:obj} 把打折清仓纳入收益口径，与商超
"运损和品相变差商品打折销售"的实际经营方式一致。

\subsubsection{5.4.2 约束条件}

\begin{enumerate}[label=（\arabic*）]
\item 补货量上限：考虑损耗后，可售商品量应能覆盖预测销量，即
      $x_{c,t}(1-\bar\ell_c)\ge\hat B_{c,t}$，等价于
      $x_{c,t}\le\hat B_{c,t}/(1-\bar\ell_c)$；
\item 价格区间：$P_{c,t}$ 限定在品类历史价格区间 $[P_c^{\min},P_c^{\max}]$ 内；
\item 空间约束：$\sum_{c=1}^{6}x_{c,t}\le V$，其中日展示容量 $V$ 取历史 6--7 月
      日总销量的 95\% 分位（$V=565.5$ kg，另取下限 $27\times2.5=67.5$ kg 保证可行）；
\item 非负约束：$x_{c,t}\ge0$。
\end{enumerate}

\subsubsection{5.4.3 方案A：成本加成定价下的线性规划}

方案 A 贴合商超的经营惯例，定价不由优化器决定：以历史平均加成率 $\eta_c$ 对预测
批发价加成，再截断至历史价格区间，即

\begin{equation}
P_{c,t}=\min\left(\max\left((1+\eta_c)\hat w_{c,t},\ P_c^{\min}\right),\ P_c^{\max}\right).
\label{eq:markup}
\end{equation}

六品类的历史平均加成率分别为 0.669、0.550、0.554、0.686、0.703 与 0.551。定价
确定后，式\eqref{eq:obj} 中目标函数只含补货量 $x_{c,t}$ 一个决策变量，约束均为
线性，模型退化为线性规划，可用单纯形类算法精确求解。

\subsubsection{5.4.4 方案B：价格离散化下的 0-1 混合整数规划}

方案 B 把定价也纳入优化。将各品类价格区间等分为 $K=10$ 档候选价 $P_c^{(k)}$，
引入 0-1 选价变量 $z_{c,t}^{(k)}$，每天每品类恰选一档，补货量不得超过所选档位
对应的销量（由三次关系式\eqref{eq:cubic}给出）：

\begin{equation}
\sum_{k=1}^{K}z_{c,t}^{(k)}=1,\qquad
x_{c,t}\le\sum_{k=1}^{K}B_c(P_c^{(k)})\,z_{c,t}^{(k)},\qquad z_{c,t}^{(k)}\in\{0,1\}.
\label{eq:milp}
\end{equation}

代入式\eqref{eq:obj} 得到混合整数规划，求解后各品类的定价与补货量联合确定。
方案 B 的可行域包含方案 A 的定价，故其最优利润不低于方案 A；两个方案的差距度量了
成本加成惯例相对定价优化的利润空间。

\subsection{5.5 求解与结果}

模型用 MATLAB 求解，方案 A 与方案 B 的求解时间均在秒级。方案 A 的逐日补货量与
定价见表4、表5，七日总利润 6,735.03 元；方案 B 的七日总利润 25,960.04 元，约为
方案 A 的 3.9 倍。图11 给出两个方案的逐日利润对比。

\begin{table}[H]
\centering
\caption{方案A 未来一周各品类日补货量（kg）}
\label{tab:q2xA}
\small
\begin{tabular}{lccccccc}
\toprule
品类 & 7月1日 & 7月2日 & 7月3日 & 7月4日 & 7月5日 & 7月6日 & 7月7日 \\
\midrule
花叶类 & 234.3 & 237.4 & 161.3 & 162.2 & 166.4 & 169.8 & 179.1 \\
花菜类 & 55.8 & 55.1 & 40.4 & 42.7 & 41.5 & 38.1 & 37.8 \\
水生根茎类 & 31.2 & 36.7 & 20.2 & 21.5 & 21.5 & 20.1 & 20.8 \\
茄类 & 41.7 & 44.3 & 28.9 & 26.3 & 27.2 & 28.5 & 29.1 \\
辣椒类 & 94.1 & 89.9 & 61.3 & 59.8 & 67.5 & 67.6 & 65.5 \\
食用菌 & 64.7 & 64.3 & 43.9 & 40.7 & 45.1 & 44.3 & 48.3 \\
\bottomrule
\end{tabular}
\end{table}

\begin{table}[H]
\centering
\caption{方案A 未来一周各品类定价（元/kg）}
\label{tab:q2pA}
\small
\begin{tabular}{lccccccc}
\toprule
品类 & 7月1日 & 7月2日 & 7月3日 & 7月4日 & 7月5日 & 7月6日 & 7月7日 \\
\midrule
花叶类 & 5.79 & 5.71 & 5.61 & 5.60 & 5.62 & 5.62 & 5.59 \\
花菜类 & 9.67 & 9.62 & 9.54 & 9.54 & 9.73 & 9.73 & 9.69 \\
水生根茎类 & 10.08 & 10.22 & 10.20 & 10.30 & 10.54 & 10.36 & 10.35 \\
茄类 & 9.44 & 9.62 & 9.46 & 9.43 & 9.51 & 9.44 & 9.45 \\
辣椒类 & 9.55 & 9.44 & 9.23 & 9.43 & 9.61 & 9.45 & 9.45 \\
食用菌 & 8.38 & 8.35 & 8.35 & 8.55 & 8.40 & 8.33 & 8.42 \\
\bottomrule
\end{tabular}
\end{table}

\begin{figure}[H]
\centering
\includegraphics[width=0.62\textwidth]{figures/q2_profit_compare.png}
\caption{方案A与方案B 日利润对比}
\label{fig:q2_profit}
\end{figure}

为检验策略的有效性，以方案 A 的口径计算 2020 至 2023 年 7 月 1 日至 7 日的利润，
结果见图12。2023 年预测年的逐日利润均高于 2021 年与 2022 年同期，处于 824 至
1,243 元区间；2020 年同期利润偏高（7 月初的 1,585 元），与该年疫情解封后消费
回补的时点有关。考虑到 2021、2022 年蔬菜零售受外部环境影响，本方案相对近两年的
提升说明补货与定价策略有效，且该对比可用实际经营背景解释。

\begin{figure}[H]
\centering
\includegraphics[width=0.68\textwidth]{figures/q2_yoy_profit.png}
\caption{2020--2023 年 7 月 1--7 日单日利润同比（方案A 口径）}
\label{fig:q2_yoy}
\end{figure}

\subsection{5.6 灵敏度分析}

对方案 A 的定价整体缩放 90\% 至 110\%，重解线性规划，结果见表6。定价每变动 5\%，
七日利润约同向变动 16\%，说明定价是利润的第一杠杆，成本加成率的设定对收益影响
显著，商超应结合竞争环境审慎调整加成水平。

\begin{table}[H]
\centering
\caption{定价灵敏度分析（方案A 口径）}
\label{tab:sensP}
\begin{tabular}{lccccc}
\toprule
定价相对基准 & 90\% & 95\% & 100\% & 105\% & 110\% \\
\midrule
七日总利润（千元） & 4.58 & 5.66 & 6.74 & 7.81 & 8.89 \\
\bottomrule
\end{tabular}
\end{table}

对空间容量 $V$ 在 0.8 至 1.2 倍区间扫描，0.8 倍容量时利润下降 3.9\%（由 6,735 元
降至 6,475 元），容量不低于基准值时利润不变，说明基准容量下空间约束不紧，模型对
容量假设稳健。
